%% summary.tex
%% Condensed technical-note summary of the HunStat2 / AD5941_25 project audit,
%% firmware fix, and dummy-cell validation (journal/01, 02, 03).
%% Based on the bare_jrnl.tex IEEEtran journal skeleton (Michael Shell, V1.4b, 2015/08/26)
%% Compile with: pdflatex summary.tex (run twice), IEEEtran.cls required.

\documentclass[journal]{IEEEtran}

\usepackage{graphicx}
\usepackage{array}
\interdisplaylinepenalty=2500

\begin{document}

\title{Project Summary: Audit, Firmware Fix, and Dummy-Cell Validation\\ of the HunStat2 Object-Oriented Potentiostat}

% [DATA REQUIRED]: real author names/affiliations, see manuscript.tex.
\author{Author~Name$^{1}$, Author~Name$^{2}$
\thanks{$^{1,2}$Authors are with [DATA REQUIRED: institution / department]. This is a condensed companion summary to the full manuscript, \texttt{manuscript.tex}, in the same directory.}}

\maketitle

\begin{abstract}
This technical note summarizes a three-part engineering review of the HunStat2 potentiostat project (AD5941 AFE, Seeed XIAO RP2040): (1) a project and object-oriented (OOP) architecture audit, (2) diagnosis and correction of two firmware defects responsible for previously-reported unusable chronoamperometry (CA) and square-wave voltammetry (SWV) data, and (3) hardware validation of the corrected firmware against a resistive/capacitive dummy cell. It is intended as a quick-reference companion to the full manuscript and the underlying \texttt{journal/01\_project\_audit.md}, \texttt{02\_oop\_analysis.md}, and \texttt{03\_firmware\_fix\_and\_dummy\_cell\_validation.md} reports.
\end{abstract}

\begin{IEEEkeywords}
Potentiostat, AD5940, firmware audit, object-oriented design, dummy-cell validation.
\end{IEEEkeywords}

\section{Project Audit Summary}
The active firmware (\texttt{Software/update~4/AD5941\_25}) targets an AD5941 electrochemical AFE controlled by a Seeed XIAO RP2040 over bit-banged SPI, with a Python/Tkinter host GUI over USB-CDC serial at 1{,}000{,}000~baud. Six electrochemical methods are implemented (OCP, CV, CA, SWV, DPV, EIS), split between a partially-modularized \texttt{src/} class tree and legacy top-level files (CV's engine was never migrated into \texttt{src/}). The repository already contained substantial prior documentation (\texttt{laporan/*.md}) and a KiCad PCB project (\texttt{STEI\_Bstat\_v1.2}); a hardware/schematic discrepancy (MCU labeled ``XIAO ESP32'' on one schematic vs.\ RP2040 everywhere else, including the actual compiled firmware target) was flagged as needing resolution before any paper states a board name. Full detail: \texttt{01\_project\_audit.md}.

\section{OOP Architecture Audit Summary}
Verdict: \textbf{partially object-oriented}, at the weak end. Eight classes exist; only two (\texttt{C\_Communication}, \texttt{C\_AD5941\_Setup}) have genuine encapsulated private state. The central parameter class, \texttt{C\_DataStorage}, has entirely public fields. No inheritance, virtual functions, or interfaces exist anywhere in the project's own code. The five electrochemical-method classes were, prior to this work, copy-paste siblings rather than a family sharing a common interface, and a parallel free-function implementation of CA/SWV/DPV plus command processing exists as confirmed dead code. Full detail and a per-file table: \texttt{02\_oop\_analysis.md}.

\section{Firmware Fix and Validation Summary}
Two defects were found and corrected in \texttt{C\_CA}, \texttt{C\_SWV}, \texttt{C\_DPV} by comparison against Analog Devices' own AD5940 example code:
\begin{enumerate}
\item \textbf{Wrong signal path}: the High-Speed (HSDAC/HSTIA) loop was configured where the Low-Power (LPDAC/LPTIA) loop --- used by every ADI DC-method reference example --- was required.
\item \textbf{Wrong ADC zero-point}: raw ADC codes were cast to \texttt{int16\_t} (zero-point at \texttt{0x0000}) instead of offset from true midscale (\texttt{0x8000}), turning small real signals into large, discontinuous, non-physical values.
\end{enumerate}
Both were fixed and validated live on hardware (COM8, 3-electrode dummy cell). Post-fix: CA shows a correctly-signed, low-noise capacitive transient on real voltage steps; CV (unmodified path) shows a near-ideal linear response, $R^2=0.999997$; SWV and DPV show smooth, bounded differential-current curves with no railing. Full detail, evidence, and data: \texttt{03\_firmware\_fix\_and\_dummy\_cell\_validation.md} and \texttt{AD5941\_25/hasil\_fixed/}.

\section{Status and Next Steps}
\begin{itemize}
\item \textbf{Done}: audit, OOP analysis, CA/SWV/DPV signal-chain fix, hardware validation of CA/SWV/DPV/CV/OCP.
\item \textbf{Open}: EIS not re-tested; dead duplicate code (\texttt{electrochemical\_methods.cpp}, \texttt{command\_processing.cpp}) not yet removed; MCU schematic discrepancy unresolved; full resolution/accuracy/noise/repeatability characterization not yet performed; CV path's absolute current-unit calibration not independently re-derived.
\item \textbf{Recommended before submission of the full manuscript}: resolve all \textbf{[DATA REQUIRED]}/\textbf{[REFERENCE REQUIRED]}/\textbf{[CODE VERIFICATION REQUIRED]} markers in \texttt{manuscript.tex} (author list, dummy-cell component values, formal citations, CV unit calibration).
\end{itemize}

\end{document}
